% Claim statement file for row claim_3_23 -- "SigmaOpenPathsWalks".
% Auto-generated stub by scaffold/solve_chapter.py at row start. The
% manager agent (or a worker it dispatches) fills the body below with the
% claim's statement(s) from the lecture notes -- statement only, no proof
% (proofs live in the sibling `claim_3_23_proof_*.tex` file).
%
% This file is a sibling of `main.tex` inside `<subsection>/tex/`. The
% `subfiles` package lets it render both standalone and as part of
% `main.tex`. Note: `main.tex` does NOT \subfile this statement file -- the
% sibling `_proof_` file restates the statement above the proof, so
% including both in `main.tex` would be redundant. This file exists for
% standalone reading / grep convenience.

\documentclass[main]{subfiles}

\begin{document}
% ref: claim_3_23 (statement)
% title: SigmaOpenPathsWalks
\def\rowref{claim\_3\_23}
\phantomsection\label{claim_3_23}

\begin{Note}[SigmaOpenPathsWalks]{Prp}{restateprpsigmaopens}\label{prp:sigma_opens}
    Let $G = (J, V, E, L)$ be a CDMG (def \ref{def-cdmg}); throughout we use the notation of def \ref{not-cdmg}, the walk and path definitions of def \ref{def:walks}, items~i.\ (a \emph{walk} in $G$ is a finite alternating sequence $\pi = (v_0, a_0, v_1, a_1, \dots, v_{n-1}, a_{n-1}, v_n)$ for some integer $n \ge 0$, with $v_0, \dots, v_n \in J \cup V$ and $a_0, \dots, a_{n-1} \in E \cup L$, satisfying the walk constraint at every $k \in \{0, 1, \dots, n - 1\}$:
    \[
        \bigl(a_k = (v_k, v_{k+1}) \;\land\; a_k \in E \cup L\bigr)
        \;\;\lor\;\;
        \bigl(a_k = (v_{k+1}, v_k) \;\land\; a_k \in E\bigr),
    \]
    and ``the same node may appear multiple times in a walk'') and~v.\ (a \emph{path} in $G$ is a walk in which no node occurs more than once), the collider/non-collider classification of def \ref{def:collider_noncollider}, item~ii.\ (collider position on a walk), and the $C$-$\sigma$-open predicate of def \ref{def:sigma_blocking}, item~1.\ (``$\pi$ is $C$-$\sigma$-open'', with $\Anc^G(C) := \bigcup_{c \in C} \Anc^G(c)$ from def \ref{def_3_5}, item~iv.). Let $C \subseteq J \cup V$ be a subset of nodes ($C$ is \emph{not} assumed disjoint from $J$, nor non-empty), and let $w_1, w_2 \in J \cup V$ be two nodes ($w_1$ and $w_2$ are \emph{not} assumed distinct: the corner case $w_1 = w_2$ is admissible and witnessed across all three statements below by the length-$0$ trivial walk $\pi = (v_0)$ with $v_0 = w_1 = w_2$, admitted by def \ref{def:walks}, item~i.\ at $n = 0$; it has no positions in $\{0, 1, \dots, n - 1\}$ and no interior positions $k \in \{1, \dots, n - 1\}$, hence no walk-edges, no colliders, and no blockable non-colliders, so it is vacuously a path, vacuously $C$-$\sigma$-open, and vacuously has every collider in $C$).

    Then the following three existential statements are equivalent:
    \begin{enumerate}
        \item there exist an integer $n \ge 0$ and a walk $\pi = (v_0, a_0, v_1, a_1, \dots, v_{n-1}, a_{n-1}, v_n)$ in $G$ in the sense of def \ref{def:walks}, item~i., with $v_0 = w_1$ and $v_n = w_2$, such that
            \[
                \bigl(\pi \text{ is a path in } G \text{ in the sense of def \ref{def:walks}, item~v.}\bigr) \;\land\; \bigl(\pi \text{ is } C\text{-}\sigma\text{-open in the sense of def \ref{def:sigma_blocking}, item~1.}\bigr);
            \]
            spelled out, ``$\pi$ is a path'' is $\forall\, k, l \in \{0, 1, \dots, n\},\; k \ne l \implies v_k \ne v_l$, and ``$\pi$ is $C$-$\sigma$-open'' is the predicate of def \ref{def:sigma_blocking}, item~1.\ (every collider position $v_k$ on $\pi$ lies in $\Anc^G(C)$, and every blockable non-collider position $v_k$ on $\pi$, in the sense of def \ref{def:unblockable_noncollider}, lies outside $C$);

        \item there exist an integer $n \ge 0$ and a walk $\pi = (v_0, a_0, v_1, a_1, \dots, v_{n-1}, a_{n-1}, v_n)$ in $G$ in the sense of def \ref{def:walks}, item~i., with $v_0 = w_1$ and $v_n = w_2$, such that $\pi$ is $C$-$\sigma$-open in the sense of def \ref{def:sigma_blocking}, item~1.;

        \item there exist an integer $n \ge 0$ and a walk $\pi = (v_0, a_0, v_1, a_1, \dots, v_{n-1}, a_{n-1}, v_n)$ in $G$ in the sense of def \ref{def:walks}, item~i., with $v_0 = w_1$ and $v_n = w_2$, such that
            \[
                \bigl(\pi \text{ is } C\text{-}\sigma\text{-open in the sense of def \ref{def:sigma_blocking}, item~1.}\bigr) \;\land\; \bigl(\forall\, k \in \{0, 1, \dots, n\},\; (k \text{ is a collider on } \pi) \implies v_k \in C\bigr),
            \]
            where ``$k$ is a collider on $\pi$'' is in the sense of def \ref{def:collider_noncollider}, item~ii. (The second conjunct strengthens the collider clause already enforced by $C$-$\sigma$-openness in def \ref{def:sigma_blocking}, item~1., from $v_k \in \Anc^G(C)$ to $v_k \in C$ --- this is the literal LN parenthetical ``and not just in $\Anc^G(C)$''. The non-collider clause of $C$-$\sigma$-openness is unchanged between (2) and (3).)
    \end{enumerate}
\end{Note}

\end{document}
